\documentclass[11pt]{article}
\usepackage{icmlko}
\begin{document}

\title{빠진 팔 하나 --- 가림만 걷으면 앱이 $-0.15$ 로 무너진다}
\author{월드모델 실험실}
\date{노트 612 · 2026년 8월 5일}
\maketitle

\begin{abstract}
논문 $462$ 가 이번 세션의 채택(노트 $586$)을 다시 재서 ``\textbf{판 이득이
고친 넷이 아니라 이웃에서 왔다}''를 냈다. 원장을 다시 읽으니 더 큰 것이
있었다 --- 노트 $583$ 은 네 팔을 쟀는데(A 옛 챔피언 $\cdot$ B 되돌림만
$\cdot$ C 되돌림$+$해제 $\cdot$ D 무작위 대조) \textbf{``되돌리지 않고
해제만'' 하는 팔이 없다.} 채택의 인과 주장 --- \emph{방향을 세웠기에 가림을
걷을 수 있다} --- 이 두 달 동안 미검증이었다.

\begin{center}
\begin{tabular}{lrrr}
\hline
팔 & 판 & KR 만화 & 비게임 앱 \\
\hline
A 옛 챔피언 & $+0.4679$ & $+0.6374$ & $+0.2897$ \\
C 채택(되돌림$+$해제) & $+0.4694$ & $+0.6765$ & $\mathbf{+0.5084}$ \\
\textbf{E 해제만} & $+0.4679$ & $+0.6626$ & $\mathbf{+0.1402}$ \\
\hline
\end{tabular}
\end{center}

\textbf{가림만 걷으면 앱이 $-0.1495$ 만큼 무너진다}
($[-0.183, -0.117]$, $0/6$). 노트 $441$ 이 \texttt{OOC\_DROP} 을 만든
근거가 정확히 재현됐고, \textbf{방향 되돌림의 순 효과는 $+0.368$} 이다.
\end{abstract}

\begin{figure}[t]
\centering
\includegraphics[width=\textwidth]{figs/f1.pdf}
\caption{\textbf{가} --- 가림만 걷으면 앱이 무너진다.
\textbf{나} --- 앱에서 부호가 뒤집힌다.
\textbf{다} --- 가림은 처방이었고 진단은 방향이었다.}
\end{figure}

\section{두 짝이 갈린다}

\begin{center}
\begin{tabular}{lrrl}
\hline
 & C $-$ A & \textbf{E $-$ A} & E 의 표본 $95\%$ \\
\hline
KR 만화 & $+0.0390$ & $+0.0252$ & $[-0.011, +0.053]$ \textbf{$0$ 을 문다} \\
비게임 앱 & $+0.2187$ & $\mathbf{-0.1495}$ & $[-0.183, -0.117]$ \textbf{$0$ 밖 음수} \\
\hline
\end{tabular}
\end{center}

\textbf{앱에서는 사전 등록 ①이 그대로 선다} --- 방향 없는 축을 안 배운
도메인에 주면 해롭다. KR 에서는 해제만으로도 $64\%$ 가 나오지만 구간이
$0$ 을 문다. \textbf{결론은 앱이 짊어진다.}

\begin{quote}
왜 앱만인가 --- 노트 $587$ 이 이미 답을 냈다. \textbf{앱은 축 무게 순서가
학습과 뒤집혀 있고 KR 은 같다.} 뒤집힌 쪽에 방향 없는 축을 주면 모형이
반대로 쓴다. \textbf{고칠 것이 있어야 고쳐진다}는 같은 문장이 여기서도
성립한다.
\end{quote}

\section{배선 가드가 정확히 $0$ 을 냈다}

E 와 A 는 \texttt{OOC\_DROP} 만 다르다. 그 손잡이는 \emph{집 밖} 도메인에만
걸리므로 판은 안 움직여야 한다.

\begin{center}
판 $E - A = \mathbf{+0.00000}$ \quad 폭 $\mathbf{0.00000}$ \quad $0/6$
\end{center}

\textbf{여섯 씨앗에서 소수점 아래가 전부 같다.} 노트 $458$ 이 세운 규칙
(``팔 둘이 정확히 같은 수를 내면 배선이 안 붙은 것이다'')의 \emph{반대}
방향 용법이다 --- \textbf{여기서는 같아야 맞다.}

\section{채택의 인과가 섰다 --- 그런데 논문 $462$ 도 틀리지 않았다}

\begin{center}
\begin{tabular}{lr}
\hline
채택이 낸 것 & 크기 \\
\hline
판(고친 넷) & $\mathbf{-0.0009}$ \\
판(이웃 팝업$\cdot$웹툰) & $+0.0031$ \\
\textbf{앱 --- 방향이 낸 몫} & $\mathbf{+0.368}$ \\
\hline
\end{tabular}
\end{center}

\textbf{둘 다 참이다.} 판에서는 방향 되돌림이 고친 도메인에 아무것도 못
줬고(오히려 $-0.0009$), 판 이득은 이웃 효과였다. \textbf{집 밖에서는 그
되돌림이 전부다} --- 없으면 $-0.15$, 있으면 $+0.22$.

\begin{quote}
\textbf{같은 변경이 자마다 다른 크기와 다른 기제로 나타난다.} 판은
\emph{배운} 도메인을 채점하므로 도메인 지시자와 함께 쓰면 그만이고 방향이
필요 없다. 집 밖 짝은 \emph{안 배운} 도메인이라 도메인 지시자가 없고,
\textbf{방향만이 그 축을 쓸 수 있게 한다.} 노트 $583$ 의 기제 설명이 두 달
만에 직접 확인됐다.
\end{quote}

\section*{조항}

\begin{quote}
\textbf{두 손잡이를 함께 채택했으면 각각만 돌린 팔이 있어야 한다.} 노트
$583$ 은 A $\cdot$ B(되돌림만) $\cdot$ C(둘 다) $\cdot$ D(무작위)를 쟀다.
\textbf{$2\times2$ 에서 한 칸이 비어 있었고}, 그 빈 칸이 하필 ``채택의
인과''를 결정하는 칸이었다. \textbf{대조 팔을 늘리는 것보다 격자를 채우는
것이 먼저다.}
\end{quote}

\section*{남은 것}

\begin{itemize}
\item \textbf{KR 이 왜 $0$ 을 무는가} --- 해제만으로 $64\%$ 가 나온다.
KR 은 축 무게 순서가 학습과 같으니(노트 $587$) 방향이 덜 필요한 것이
설명이지만, \textbf{그 설명을 KR 쪽에서 따로 재지 않았다.}
\item \textbf{셋째 짝 CN 만화를 안 넣었다} --- $352$ 행이라 표본 SD 가
크겠지만, 거울 조항(노트 $426$)은 짝 둘로만 서 있다.
\item \textbf{$B$ 팔을 표본 구간으로 다시 안 봤다} --- 노트 $583$ 의
``KR $-0.0111$ $\cdot$ 앱 $-0.0176$ 둘 다 $0$ 밖 음수''는 \textbf{씨앗
구간}이다. 노트 $608$ 이후 그 표현은 못 쓴다.
\end{itemize}

\end{document}
